\chapter{Architecture}
\label{ch:architecture}

This chapter defines the v2 system architecture: the five primary components, their interfaces and invariants,
the complete data flow from runbook invocation to trace archive, the execution lifecycle, the concurrency model,
and how the \texttt{gert serve} JSON-RPC layer integrates as a first-class adapter.

The single most important architectural rule is: \textbf{dependencies flow inward toward the Core Domain.}
The Runtime Core knows nothing about TUI, VS Code, or JSON-RPC. Adapters know nothing about execution semantics.
Every component boundary is enforced by a Go interface; no component imports a package that is ``above'' it in
the dependency graph.

% ─────────────────────────────────────────────────────────────────────────────
% Figure: Architecture dependency graph
% ─────────────────────────────────────────────────────────────────────────────
\begin{figure}[htbp]
\centering
\begin{tikzpicture}[node distance=1.0cm and 1.5cm, scale=0.9, transform shape]
  %% ── Pipeline (top to bottom) ────────────────────────────────────────────────
  \node[stepbox]                          (arc-runbook)   {Runbook (YAML)};
  \node[stepbox, below=of arc-runbook]    (arc-parser)    {Parser};
  \node[stepbox, below=of arc-parser]     (arc-planner)   {Planner};
  \node[stepbox, below=1.5cm of arc-planner] (arc-executor) {Executor};

  %% ── RPC Server and Audit Log (right side) ───────────────────────────────────
  \node[stepbox, right=2.5cm of arc-executor]      (arc-rpcserver) {RPC Server};
  \node[stepbox, below=0.5cm of arc-rpcserver]     (arc-auditlog)  {Audit Log (JSONL)};

  %% ── Runtime Core sub-components ─────────────────────────────────────────────
  \node[stepbox, below left=1.3cm  and 0.0cm of arc-executor]  (arc-handlers) {Step Handlers};
  \node[stepbox, below right=1.3cm and 0.0cm of arc-executor]  (arc-varstore) {Variable Store};

  %% ── Services (driven by Step Handlers) ──────────────────────────────────────
  \node[stepbox] at ([xshift=-3.5cm, yshift=-3.7cm] arc-executor) (arc-toolreg)  {Tool Registry};
  \node[stepbox] at ([xshift=-0.5cm, yshift=-3.7cm] arc-executor) (arc-inputprov){Input Provider};
  \node[stepbox] at ([xshift=+2.5cm, yshift=-3.7cm] arc-executor) (arc-eventbus) {Event Bus};

  %% ── Runtime Core grouping box ────────────────────────────────────────────────
  \begin{scope}[on background layer]
    \node[gertfit, fit=(arc-executor)(arc-handlers)(arc-varstore)] (arc-rtcore) {};
  \end{scope}
  \node[below=3pt of arc-rtcore, font=\small\sffamily\itshape, color=black!50]
       {Runtime Core};

  %% ── Arrows ───────────────────────────────────────────────────────────────────
  \draw[gertarrow] (arc-runbook)   -- (arc-parser);
  \draw[gertarrow] (arc-parser)    -- (arc-planner);
  \draw[gertarrow] (arc-planner)   -- (arc-executor);
  \draw[gertarrow] (arc-rpcserver) -- (arc-executor);
  \draw[gertarrow] (arc-executor)  -- (arc-auditlog);
  \draw[gertarrow] (arc-executor)  -- (arc-handlers);
  \draw[gertarrow] (arc-executor)  -- (arc-varstore);
  \draw[gertarrow] (arc-handlers)  -- (arc-toolreg);
  \draw[gertarrow] (arc-handlers)  -- (arc-inputprov);
  \draw[gertarrow] (arc-handlers)  -- (arc-eventbus);
\end{tikzpicture}
\caption{gert v2 architecture --- component dependency overview.}
\label{fig:architecture-overview}
\end{figure}

% ─────────────────────────────────────────────────────────────────────────────
% Figure: Dependency direction convention legend
% ─────────────────────────────────────────────────────────────────────────────
\begin{figure}[H]
\centering
\begin{tikzpicture}[node distance=0.6cm and 1.8cm, scale=0.9, transform shape]
  %% ── Example 1: A depends on B ───────────────────────────────────────────────
  \node[stepbox] (dd-a1) {A};
  \node[stepbox, right=of dd-a1] (dd-b1) {B};
  \draw[gertarrow] (dd-a1) -- (dd-b1);
  \node[below=0.25cm of dd-a1, font=\small\itshape, color=black!60] {caller};
  \node[below=0.25cm of dd-b1, font=\small\itshape, color=black!60] {dependency};
  \node[font=\small, right=0.4cm of dd-b1] (dd-label1)
       {\textrightarrow\ \textit{``A depends on B''}};

  %% ── Example 2: Transitive dependency ────────────────────────────────────────
  \node[stepbox, below=1.4cm of dd-a1] (dd-a2) {A};
  \node[stepbox, right=of dd-a2]       (dd-b2) {B};
  \node[stepbox, right=of dd-b2]       (dd-c2) {C};
  \draw[gertarrow] (dd-a2) -- (dd-b2);
  \draw[gertarrow] (dd-b2) -- (dd-c2);
  \draw[gertarrow, dashed, bend right=35] (dd-a2) to
       node[below, font=\scriptsize\itshape, color=black!55] {transitive} (dd-c2);
  \node[font=\small, right=0.4cm of dd-c2] (dd-label2)
       {\textrightarrow\ \textit{``A transitively depends on C''}};
\end{tikzpicture}
\caption{Dependency direction convention: an arrow from A to B means \emph{A depends on B} (A imports/uses B). Arrows flow \emph{toward} the depended-upon component.}
\label{fig:dependency-direction}
\end{figure}

% ─────────────────────────────────────────────────────────────────────────────
\section{Primary Components}
% ─────────────────────────────────────────────────────────────────────────────

\subsection{Parser}

\paragraph{Responsibility.}
The Parser owns the transformation of a raw YAML document into a validated, semantically-resolved
\texttt{ParsedRunbook} value. It does \emph{not} own execution, planning, tool discovery, or input resolution.
It validates schema structure and surface-level semantic rules (duplicate step IDs, unknown step types,
malformed Go template expressions) but intentionally defers deep semantic validation (e.g., whether a
\texttt{from:} binding can actually be resolved) to later pipeline stages.

\paragraph{Interface.}

\begin{minted}{go}
// Parser transforms raw YAML bytes into a validated ParsedRunbook.
type Parser interface {
    // Parse reads YAML from r, validates against the v2 JSON Schema,
    // applies structural semantic checks, and returns a ParsedRunbook.
    // Errors carry structured diagnostics (file, line, column, message).
    Parse(ctx context.Context, r io.Reader, opts ParseOptions) (*ParsedRunbook, error)
}

type ParseOptions struct {
    // SchemaVersion overrides the apiVersion field for compatibility shims.
    SchemaVersion string
    // StrictMode causes deprecation warnings to be surfaced as errors.
    StrictMode bool
}

// ParsedRunbook is an immutable, fully-validated representation of a runbook.
// All fields have been structurally validated. Templates are not yet evaluated.
type ParsedRunbook struct {
    Meta       RunbookMeta
    Governance GovernancePolicy  // nil if no governance block
    Steps      []ParsedStep
    Imports    map[string]string // alias -> path (not yet resolved)
    Tools      []string          // tool names (not yet discovered)
    Schema     SchemaRef         // version and $schema URL
}
\end{minted}

\paragraph{Dependencies.}
The Parser depends only on the JSON Schema artifact (a static embedded file) and the expression
evaluator (for template syntax validation). It has no runtime dependencies.

\paragraph{Key Invariants.}
\begin{itemize}
  \item A \texttt{ParsedRunbook} returned without error is structurally valid against the v2 JSON Schema.
  \item All step IDs are unique within a \texttt{ParsedRunbook}.
  \item No I/O occurs during parsing (file imports are recorded as paths, not resolved).
  \item Parsing is pure and deterministic: the same input always produces the same output.
\end{itemize}

% ──────────────────────────────────────────────────────────────────────────────
\subsection{Planner}

\paragraph{Responsibility.}
The Planner transforms a \texttt{ParsedRunbook} into an \texttt{ExecutionPlan}: a resolved, ordered,
dependency-annotated representation of every step the runtime will execute. The Planner resolves imports
(recursively parsing and inlining included runbooks), discovers and validates tool definitions, resolves
input provider bindings to their declared sources, and produces the static step ordering. It does
\emph{not} evaluate dynamic conditions (branch predicates, iterate count expressions) — those are
runtime concerns.

\paragraph{What is an ExecutionPlan?}
An \texttt{ExecutionPlan} is a \emph{flat, ordered list} of resolved steps with pre-computed metadata,
not a DAG. Branch logic and iteration are left as runtime decisions because they depend on captured
variables evaluated during execution. The plan encodes \emph{which steps exist} and \emph{what
preconditions they carry}, not which steps will actually run. This design keeps the planner pure and
the runtime authoritative over control flow.

\begin{minted}{go}
type Planner interface {
    // Plan resolves imports, discovers tools, validates bindings,
    // and returns a fully-resolved ExecutionPlan ready for the runtime.
    Plan(ctx context.Context, rb *ParsedRunbook, opts PlanOptions) (*ExecutionPlan, error)
}

type PlanOptions struct {
    BaseDir     string            // working directory for import resolution
    ToolDirs    []string          // directories to search for .tool.yaml files
    ProviderDir string            // directory for .provider.yaml files
    Vars        map[string]string // CLI-supplied variable overrides
}

// ExecutionPlan is an immutable, fully-resolved plan for a single run.
type ExecutionPlan struct {
    RunID       string
    RunbookPath string
    Steps       []ResolvedStep       // ordered; includes inlined include steps
    Tools       map[string]*ToolDef  // name -> resolved tool definition
    Providers   map[string]*ProviderDef // binding name -> provider
    Governance  *GovernancePolicy
    Metadata    PlanMetadata
}

// ResolvedStep is a single step in the plan, fully resolved.
type ResolvedStep struct {
    ID       string
    Kind     StepKind      // cli, choice, decision, collector, tool, include, branch, iterate
    Spec     StepSpec      // kind-specific parameters (resolved, not templated)
    Contract *contract.Contract  // effects/reads/writes/idempotent/deterministic
    Depth    int           // chain depth (0 = root runbook)
    Origin   string        // source runbook path (for inlined include steps)
}
\end{minted}

\paragraph{Dependencies.}
The Planner depends on the Parser (for import resolution), the Tool Registry (for
\texttt{.tool.yaml} discovery), and the Provider Registry (for \texttt{.provider.yaml}
discovery). It must not depend on the Runtime Core, Extension Host, or any adapter.

\paragraph{Key Invariants.}
\begin{itemize}
  \item An \texttt{ExecutionPlan} returned without error contains only resolvable steps.
  \item All tool names referenced by steps appear in \texttt{ExecutionPlan.Tools}.
  \item Import graphs are acyclic; the Planner detects and rejects cycles.
  \item All \texttt{from:} bindings that declare a provider appear in \texttt{ExecutionPlan.Providers}.
  \item The Planner does not execute or side-effect; it is safe to plan without running.
\end{itemize}

% ──────────────────────────────────────────────────────────────────────────────
\subsection{Runtime Core}

\paragraph{Responsibility.}
The Runtime Core owns the execution loop: advancing through an \texttt{ExecutionPlan} step by step,
evaluating templates and branch predicates, dispatching tool invocations, collecting evidence, writing
the append-only trace, enforcing governance pre-flight checks, emitting structured events, and
managing the mutable run state. It does \emph{not} own user presentation, serialization formats, or
network transport.

\paragraph{Interface.}

\begin{minted}{go}
type Runtime interface {
    // Start initializes a run from a plan and returns a RunHandle.
    // The run does not advance until Next is called.
    Start(ctx context.Context, plan *ExecutionPlan, opts RunOptions) (RunHandle, error)

    // Resume restores a previously checkpointed run from the run store.
    Resume(ctx context.Context, runID string, opts RunOptions) (RunHandle, error)
}

type RunHandle interface {
    // Next advances the run by one step. Blocks until the step completes
    // (or is paused waiting for evidence/approval).
    // Returns io.EOF when the run has no more steps.
    Next(ctx context.Context) (*StepResult, error)

    // Approve records approval for the current step's approval gate.
    Approve(ctx context.Context, decision ApprovalDecision) error

    // SubmitEvidence records evidence for interactive steps (choice, decision, collector).
    // For choice: stores selected option. For decision: stores route selection.
    // For collector: stores multi-field form data and artifact attachments.
    SubmitEvidence(ctx context.Context, stepID string, ev map[string]*EvidenceValue) error

    // Cancel requests a graceful stop of the run.
    Cancel(ctx context.Context, reason string) error

    // State returns the current mutable run state (read-only snapshot).
    State() RunState

    // Events returns a channel of structured events emitted during execution.
    Events() <-chan Event
}

type RunOptions struct {
    Mode        RunMode           // real, dry-run, replay
    Actor       string            // --as identity
    Vars        map[string]string // variable overrides
    ScenarioDir string            // replay scenario directory
    Store       RunStore          // durable store for checkpointing
    OnEvent     func(Event)       // synchronous event hook (nil = use Events channel)
}
\end{minted}

\paragraph{Dependencies.}
The Runtime Core depends on the Governance Engine, Tool Runtime (for tool invocation),
Expression Evaluator (for templates and branch predicates), Evidence store, and Trace Writer.
It does not depend on the Extension Host, API/Adapter Layer, or any adapter.

\paragraph{Key Invariants.}
\begin{itemize}
  \item The trace is append-only: no event is ever deleted or modified after being written.
  \item Governance pre-flight is enforced before every \texttt{cli} and \texttt{tool} step dispatch.
  \item Redaction is applied to all captured output before it is written to the trace.
  \item The run state is checkpointed after every step completion.
  \item Events are emitted in a total order that matches step execution order.
  \item The runtime is re-entrant per run: multiple callers may call \texttt{Next} sequentially
        (but not concurrently) on the same \texttt{RunHandle}.
\end{itemize}

% ──────────────────────────────────────────────────────────────────────────────
\subsection{Extension Host}

\paragraph{Responsibility.}
The Extension Host manages the lifecycle of out-of-process extensions: discovery, launch, capability
negotiation, RPC dispatch, sandboxing, and graceful shutdown. It exposes extension-contributed tools
and providers to the Tool Runtime and Provider Registry. It does \emph{not} execute runbook steps
directly and does not own the extension communication protocol (that belongs to the extension RPC
layer, defined in §04).

\paragraph{Interface.}

\begin{minted}{go}
type ExtensionHost interface {
    // Load discovers and launches all extensions declared in the project manifest.
    Load(ctx context.Context, manifest *ProjectManifest) error

    // Shutdown gracefully stops all running extensions.
    Shutdown(ctx context.Context) error

    // ContributedTools returns all tools contributed by loaded extensions.
    ContributedTools() []*ToolDef

    // ContributedProviders returns all providers contributed by loaded extensions.
    ContributedProviders() []*ProviderDef

    // ContributedPolicyRules returns governance rules contributed by extensions.
    // These compose with host policy; host policy takes precedence on conflict.
    ContributedPolicyRules() []PolicyRule
}
\end{minted}

\paragraph{Dependencies.}
The Extension Host depends on the extension manifest schema and the extension RPC transport
(stdio JSON-RPC, as specified in §04). It does not depend on the Runtime Core or any adapter.

\paragraph{Key Invariants.}
\begin{itemize}
  \item No extension operates without an explicit capability grant from the host.
  \item Extension crashes do not crash the host process.
  \item The Extension Host enforces timeout and cancellation on all RPC calls to extensions.
  \item Extension-contributed policy rules are additive; they can only \emph{restrict}, not
        \emph{relax}, host-defined policy.
\end{itemize}

% ──────────────────────────────────────────────────────────────────────────────
\subsection{API / Adapter Layer}

\paragraph{Responsibility.}
The API/Adapter Layer is the boundary between the Runtime Core and all external consumers
(TUI, VS Code extension, web clients, CI pipelines). It translates the runtime's event stream
and \texttt{RunHandle} contract into transport-specific protocols. Each adapter is a thin,
stateless renderer: it subscribes to the event stream, presents state to the user, and forwards
user input (approvals, evidence, cancellation) back to the \texttt{RunHandle}. Adapters own
\emph{no} execution logic.

\paragraph{Interface.}

\begin{minted}{go}
// Adapter is the interface every consumer of the Runtime Core must implement.
type Adapter interface {
    // Attach binds the adapter to a RunHandle before the first Next call.
    Attach(handle RunHandle) error

    // Run blocks until the run completes or is cancelled. It is responsible
    // for draining Events(), rendering output, and forwarding user input.
    Run(ctx context.Context) error
}

// The gert serve JSON-RPC layer implements Adapter over stdio.
// The TUI (Bubble Tea) implements Adapter over a terminal.
// Future: a web adapter implements Adapter over WebSocket.
\end{minted}

\paragraph{Dependencies.}
Adapters depend on the \texttt{RunHandle} and \texttt{Event} types from the Runtime Core.
They must not import any package from the Parser, Planner, or Extension Host.

\paragraph{Key Invariants.}
\begin{itemize}
  \item An adapter may not call \texttt{Next} more than once concurrently.
  \item An adapter may not modify run state directly; all state changes go through \texttt{RunHandle}.
  \item Adapters must handle the case where the event channel is closed (run ended) without blocking.
  \item An adapter failure must not corrupt the run trace or run state.
\end{itemize}

\paragraph{Dependency Direction Summary.}

\begin{figure}[H]
\centering
\begin{tikzpicture}[node distance=1.0cm and 1.5cm]
  %% ── Vertical spine (Adapter → Runtime Core → Planner → Parser → Schema) ─────
  \node[stepbox]                          (dd-adapter) {Adapter Layer};
  \node[stepbox, below=of dd-adapter]     (dd-rtcore)  {Runtime Core};
  \node[stepbox, below=of dd-rtcore]      (dd-planner) {Planner};
  \node[stepbox, below=of dd-planner]     (dd-parser)  {Parser};
  \node[stepbox, below=of dd-parser]      (dd-schema)  {Schema / JSON Schema};

  %% ── Peer column (Extension Host + Tool Runtime, right of spine) ──────────────
  \node[stepbox, right=2.5cm of dd-rtcore]  (dd-exthost) {Extension Host};
  \node[stepbox, below=of dd-exthost]       (dd-toolrt)  {Tool Runtime};

  %% ── Arrows ───────────────────────────────────────────────────────────────────
  \draw[gertarrow] (dd-adapter)  -- (dd-rtcore);
  %% Extension Host <-> Runtime Core: peer relationship, both arrows shown
  \draw[gertarrow] (dd-exthost)  to[bend right=12] (dd-rtcore);
  \draw[gertarrow] (dd-rtcore)   to[bend right=12] (dd-exthost);
  \draw[gertarrow] (dd-exthost)  -- (dd-toolrt);
  \draw[gertarrow] (dd-rtcore)   -- (dd-planner);
  \draw[gertarrow] (dd-planner)  -- (dd-parser);
  \draw[gertarrow] (dd-parser)   -- (dd-schema);
\end{tikzpicture}
\caption{Component dependency direction. Arrows point from dependent to dependency. No upward arrows exist.}
\label{fig:dep-direction}
\end{figure}

No upward arrow exists. The Parser knows nothing about the Runtime. The Runtime knows nothing
about adapters. The Extension Host is a peer of the Runtime Core, contributing tools and policy
rules but not owning execution.

% ─────────────────────────────────────────────────────────────────────────────
\section{Data Flow}
% ─────────────────────────────────────────────────────────────────────────────

This section traces the complete path from runbook invocation to trace archive, identifying
exactly where each responsibility is exercised.

\subsection{Phase 1: Parse}

The user invokes \texttt{gert exec runbook.yaml --as eng@corp.com}. The CLI layer reads the
YAML file and calls \texttt{Parser.Parse()}. The Parser validates the document against the
embedded v2 JSON Schema, checks structural semantic rules, and returns a \texttt{ParsedRunbook}.
If validation fails, structured diagnostics are returned and execution does not proceed.

\subsection{Phase 2: Plan}

The \texttt{ParsedRunbook} is passed to \texttt{Planner.Plan()}. The Planner:
\begin{enumerate}
  \item Resolves all \texttt{imports:} recursively (parse → plan each imported runbook, inline
        its steps with incremented \texttt{Depth}).
  \item Discovers \texttt{.tool.yaml} files for each named tool.
  \item Resolves \texttt{from:} input bindings to their declared providers.
  \item Validates that all step references (branch targets, iterate bodies) are valid step IDs.
  \item Constructs and returns an \texttt{ExecutionPlan}.
\end{enumerate}

Governance policy from the runbook's \texttt{meta.governance} block is extracted and embedded
in the \texttt{ExecutionPlan} for runtime enforcement.

\subsection{Phase 3: Run Initialization}

\texttt{Runtime.Start(plan, opts)} is called. The runtime:
\begin{enumerate}
  \item Generates a \texttt{RunID} (format: \texttt{YYYYMMDDTHHmmss-xxxxxxxx}).
  \item Creates the run directory: \texttt{.runbook/runs/<runID>/}.
  \item Opens the append-only trace file: \texttt{.runbook/runs/<runID>/trace.jsonl}.
  \item Writes the \texttt{run/started} trace event (runbook path, actor, mode, plan hash).
  \item Initializes the \texttt{GovernanceEngine} from the plan's governance policy.
  \item Compiles redaction rules from the governance policy.
  \item Returns a \texttt{RunHandle}; the run is now paused at step 0.
\end{enumerate}

\subsection{Phase 4: Step Execution Loop}

The adapter calls \texttt{RunHandle.Next(ctx)} in a loop until \texttt{io.EOF}.
For each step, the runtime executes the following sequence:

\subsubsection*{Step Pre-flight (all step types)}
\begin{enumerate}
  \item Write \texttt{step/started} trace event.
  \item Emit \texttt{StepStarted} runtime event to the event channel.
  \item Evaluate template expressions in the step spec against current run variables.
\end{enumerate}

\subsubsection*{Governance Pre-flight (cli and tool steps)}
\begin{enumerate}
  \item \textbf{Command check:} \texttt{GovernanceEngine.CheckCommand(argv[0])} evaluates
        the denylist (deny takes precedence) and then the allowlist. A violation writes a
        \texttt{governance/commandBlocked} trace event and returns an error; the step fails.
  \item \textbf{Env var check:} \texttt{GovernanceEngine.FilterEnvVars()} removes variables
        matching \texttt{deny\_env\_vars} patterns. Blocked variable names are written to the
        \texttt{governance/envVarBlocked} trace event.
  \item \textbf{Contract check:} The step's \texttt{Contract.Risk()} is evaluated. Steps
        with \texttt{RiskCritical} that do not have an explicit approval gate declaration
        in the plan emit a \texttt{governance/contractWarning} event (not a hard failure in
        v2.0; upgraded to error in strict mode).
  \item \textbf{Approval gate:} If the step declares \texttt{approvals.min > 0}, the runtime
        pauses and emits an \texttt{ApprovalRequired} event. Execution blocks at this step
        until \texttt{RunHandle.Approve()} is called with sufficient approvals.
        See §\ref{sec:approval-gate} (Governance chapter) for the full approval protocol.
\end{enumerate}

\subsubsection*{Dispatch (step-type-specific)}
\begin{itemize}
  \item \textbf{cli:} Fork a subprocess. Capture stdout/stderr. Apply redaction rules to
        captured output before storing in run variables or writing to trace.
  \item \textbf{tool:} Call \texttt{ToolRuntime.Invoke(toolName, input)}. The tool runtime
        dispatches to the appropriate transport (stdio, JSON-RPC, MCP). Redaction applied
        to returned output.
  \item \textbf{choice:} Emit \texttt{ChoiceRequired} event with option list. Block until
        \texttt{SubmitEvidence()} is called with selected option. Store result as named
        variable. No artifacts involved.
  \item \textbf{decision:} Emit \texttt{DecisionRequired} event with route map. Block until
        \texttt{SubmitEvidence()} is called with route selection. Alter execution flow by
        jumping to the selected runbook or label. Optionally store selection in variable
        for audit trail.
  \item \textbf{collector:} Emit \texttt{CollectorRequired} event with field schema. Block
        until \texttt{SubmitEvidence()} is called with multi-field form data. Validate each
        submitted value against its declared field type and constraints (see
        §\ref{sec:collector_field_validation}). SHA256-hash each \texttt{file} or
        \texttt{image} attachment. Store typed values as run variables: \texttt{number} and
        \texttt{integer} as JSON numbers, \texttt{boolean} as JSON boolean, multi-select
        \texttt{select} and \texttt{choice} as JSON arrays of strings, single-select
        \texttt{select} as a JSON string, \texttt{date} and \texttt{datetime} as ISO 8601
        strings, \texttt{text} and \texttt{url} as JSON strings. Store \texttt{file} and
        \texttt{image} artifacts in the evidence store (not in run variables).
  \item \textbf{include:} Resolve the referenced runbook. Validate the include graph for
        cycles (load-time DFS — fail fast if cyclic). Evaluate \texttt{include.when} if
        present; skip all steps if false. Apply \texttt{include.with} overrides to the
        current variable scope. Inline-expand the referenced runbook's steps into the
        current execution queue at \texttt{Depth+1}. The include step itself does not
        appear in the run trace; expanded steps run as if defined inline, each carrying
        an \texttt{Origin} breadcrumb: \texttt{[included from: <path>]}.
  \item \textbf{branch:} Evaluate branch predicates (expression evaluator) against current
        variables. Select the matching branch arm. Write \texttt{event/branchResolved}
        trace event. Advance step cursor to the first step of the resolved arm.
  \item \textbf{iterate:} Evaluate the iteration expression. Write
        \texttt{event/iteratePassStart}. Execute body steps. Write
        \texttt{event/iteratePassEnd}. Repeat until condition is false.
  \item \textbf{wait\_for\_event:} Serialise run state to the run store. Register an event
        listener with the Event Dispatcher (source, event ID, filter predicate, and timeout).
        Set run state to \texttt{WAITING} and return to caller — execution is suspended.
        When the matching event arrives, the Event Dispatcher rehydrates the run, validates
        the payload against \texttt{event.payload\_schema}, applies \texttt{event.filter}
        conditions, writes captured variables via \texttt{capture} mappings, and resumes
        at the next step. If \texttt{timeout} elapses before an event arrives the
        \texttt{on\_timeout} policy governs the outcome (see
        §\ref{sec:wait_for_event_contract}). This dispatch path is only available in
        \texttt{gert serve} mode; \texttt{gert run} MUST reject the step with error:
        \texttt{"step type 'wait\_for\_event' requires gert serve mode"}.
\end{itemize}

\subsubsection*{Step Type Classification}

All step types fall into one of three categories based on runtime behavior:

\begin{center}
\begin{tabular}{|l|l|p{7cm}|}
\hline
\textbf{Category} & \textbf{Step Types} & \textbf{Runtime Behavior} \\
\hline
\textbf{Execution} & cli, tool & Execute external commands or tools. Require governance
pre-flight (allowlist, denylist, redaction). Capture output and exit codes. \\
\hline
\textbf{Interactive} & choice, decision, collector & Block execution, emit event requiring
user input. \texttt{choice} stores selected option as variable. \texttt{decision} alters
control flow by routing to a runbook or label. \texttt{collector} captures multi-field
form data and artifacts (files, images). All three write evidence to trace. \\
\hline
\textbf{Control Flow} & branch, iterate, include & Alter execution path or loop body. \texttt{branch}
evaluates predicates at runtime. \texttt{iterate} repeats a body until a condition is false.
\texttt{include} resolves and inline-expands a referenced runbook's steps into the current
execution queue; the include step itself is not emitted to the trace. \\
\hline
\textbf{Synchronisation} & wait\_for\_event & Suspend the run until an external event arrives via
webhook, message queue, OS signal, or named in-process channel. Persists full run state to durable
storage; resumes when the Event Dispatcher delivers a matching, validated event. Valid only in
\texttt{gert serve} mode — rejected at runtime in \texttt{gert run} (local CLI) mode. \\
\hline
\textbf{Approval Gate} & approve & Pause execution pending sign-off from one or more authorised
identities. Supports \texttt{all}/\texttt{any}/\texttt{quorum} modes with an explicit approver
pool. Optionally enforces a business-day timeout (GAP-1) and M-of-N quorum tracking (GAP-2). See
§\ref{sec:approve_executor_contract}. \\
\hline
\textbf{Presentation} & display & Render template content to the operator without requiring input.
\texttt{content} is evaluated via Go \texttt{text/template} against the current variable scope.
\texttt{format} controls rendering hint (\texttt{text} or \texttt{markdown}). Fire-and-continue;
no state mutation. \\
\hline
\textbf{Utility} & noop & No-op placeholder. Evaluates \texttt{capture} expressions. Zero side effects. \\
\hline
\end{tabular}
\end{center}


% ─────────────────────────────────────────────────────────────────────────────
\subsection{CLI Step Executor Contract}
\label{sec:cli_executor_contract}
% ─────────────────────────────────────────────────────────────────────────────

A \texttt{type: cli} step executes a shell command or script on the host.  Two
execution forms are supported:

\begin{itemize}
  \item \textbf{Exec form} (\texttt{command}/\texttt{args}) — bypasses the shell
        entirely.  The runtime forks the binary directly, which is both faster and
        injection-safe.
  \item \textbf{Shell-string form} (\texttt{run: "<string>"}) — passes a script to a
        shell interpreter.  Supports multi-line scripts and shell built-ins at the cost
        of increased injection risk.
  \item \textbf{Multi-shell map form} (\texttt{run: \{bash: "…", cmd: "…", …\}}) —
        provides per-shell scripts; the executor selects the first available shell at
        run time.  Designed for cross-platform runbooks.
\end{itemize}

The fields \texttt{command} and \texttt{run} are mutually exclusive.  If both are
present the step fails validation with an error (see validation table below).

\subsubsection*{Shell Resolution — String Form}

When \texttt{run} is a string the executor resolves the shell before forking.

\paragraph{Automatic detection (\texttt{shell: auto}, the default).}
The executor inspects the OS at plan-time and selects the platform-native shell:

\begin{itemize}
  \item \textbf{Linux / macOS:} \texttt{/bin/sh} — POSIX-guaranteed to exist.
  \item \textbf{Windows:} \texttt{pwsh} if \texttt{exec.LookPath("pwsh")} succeeds;
        otherwise \texttt{cmd.exe}.
\end{itemize}

\paragraph{Explicit shell (\texttt{shell: bash | sh | pwsh | cmd}).}
The executor calls \texttt{exec.LookPath(shell)} at \emph{plan time}.  If the binary
is not found the plan fails immediately with:

\begin{verbatim}
Shell 'bash' not found on this host (exec.LookPath returned empty)
\end{verbatim}

This is a plan-time error, not a run-time error — the runbook is rejected before
execution begins.

\paragraph{Invocation patterns.}
Once the shell is resolved the script is passed using the shell's standard
non-interactive flag:

\begin{center}
\begin{tabular}{|l|l|}
\hline
\textbf{Shell} & \textbf{Invocation} \\
\hline
\texttt{bash} & \texttt{["bash", "-c", script]} \\
\texttt{sh}   & \texttt{["sh", "-c", script]} \\
\texttt{pwsh} & \texttt{["pwsh", "-NonInteractive", "-Command", script]} \\
\texttt{cmd}  & \texttt{["cmd.exe", "/c", script]} \\
\hline
\end{tabular}
\end{center}

\paragraph{Template expansion.}
Template expansion (Sprig/Go template) is applied to the \texttt{run} string
\emph{before} the result is handed to the shell.  The shell receives the
already-expanded string.  Shell interpretation therefore operates on the final
expanded text, not on the template source.

\paragraph{Security note — shell injection.}
Shell-string form is \textbf{not} immune to shell injection.  If a template variable
contains shell metacharacters (\texttt{\$}, \texttt{`}, \texttt{;}, \texttt{|},
\texttt{\&}, parentheses, etc.) they will be interpreted by the shell after template
expansion.  \emph{Authors are responsible for sanitizing template values used in
shell-string steps.}  For untrusted input, prefer the exec form
(\texttt{command}/\texttt{args}), which bypasses the shell entirely.

\subsubsection*{Shell Resolution — Multi-Shell Map Form}

When \texttt{run} is a YAML mapping (e.g.\
\texttt{run: \{bash: "…", pwsh: "…"\}}) the executor performs a two-phase
resolution.

\paragraph{Phase 1 — Plan-time key validation.}
The executor reads all map keys and verifies each against the set of recognized
shell names: \texttt{bash}, \texttt{sh}, \texttt{pwsh}, \texttt{cmd}.  Any
unrecognized key is a load-time validation error:

\begin{verbatim}
Unknown shell 'zsh' in run map. Valid values: bash, sh, pwsh, cmd
\end{verbatim}

The plan is rejected before execution begins.

\paragraph{Phase 2 — Run-time shell selection.}
At the moment the step is dispatched the executor walks the map in \emph{insertion
order} and calls \texttt{exec.LookPath(key)} for each entry:

\begin{enumerate}
  \item The \textbf{first key whose shell resolves successfully} is selected.
  \item \textbf{Platform-family fallback:} if the \texttt{bash} key is present but
        \texttt{bash} is not found, the executor tries \texttt{sh} as a fallback
        before moving on to the next map key.
  \item If \textbf{no key resolves} the step fails with:
\begin{verbatim}
CLI step '<id>': no shell from map {bash, pwsh} available on this host.
Install one of: bash, pwsh.
\end{verbatim}
\end{enumerate}

The \texttt{shell} field is ignored when \texttt{run} is a map.  Setting it
produces a warning (see validation table).

\paragraph{Audit record.}
The JSONL trace records which map entry was selected.  A \texttt{shell\_selected}
field is written alongside the step's \texttt{step/started} event:

\begin{minted}{json}
{ "kind": "step/started", ..., "payload": {
    "step_id": "deploy-linux",
    "shell_selected": "bash",
    ...
  }
}
\end{minted}

This field appears in the \texttt{step/started} payload whenever a shell-string or
map-form CLI step begins.  For exec-form steps the field is omitted.

\paragraph{Template expansion.}
The same expansion-before-invocation rule applies: template expansion is applied to
the selected map value before it is passed to the shell.

\subsubsection*{Platform Variable Injection}

The runtime injects a \texttt{gert.platform} namespace into the variable scope at
run start.  These variables are available to all template expressions throughout the
run, including \texttt{branch} predicates.

\begin{center}
\begin{tabular}{|l|l|p{7cm}|}
\hline
\textbf{Variable} & \textbf{Type} & \textbf{Values} \\
\hline
\texttt{gert.platform.os}    & string & \texttt{"linux"} \textbar\ \texttt{"darwin"} \textbar\ \texttt{"windows"} \\
\texttt{gert.platform.arch}  & string & \texttt{"amd64"} \textbar\ \texttt{"arm64"} \textbar\ \texttt{"arm"} \\
\texttt{gert.platform.shell} & string & Shell name selected for the current step (e.g.\ \texttt{"bash"}).
                                        Set after shell resolution; available in \texttt{capture} output
                                        and downstream steps.  \textbf{Not} usable inside the \texttt{run}
                                        script of the same step — resolution happens before template
                                        expansion of that field is complete. \\
\hline
\end{tabular}
\end{center}

\texttt{gert.platform.os} is particularly useful for complementing the map form with
\texttt{branch} steps in cases requiring more complex platform-conditional logic than
the map form alone can express.

\subsubsection*{Validation Rules}

\begin{center}
\begin{tabular}{|p{5.5cm}|l|p{6cm}|}
\hline
\textbf{Rule} & \textbf{Severity} & \textbf{Message} \\
\hline
\texttt{command} and \texttt{run} both present & Error &
  \texttt{Fields \`{}command\`{} and \`{}run\`{} are mutually exclusive on type: cli} \\
\hline
\texttt{run} is a map \textbf{and} \texttt{shell} is set & Warning &
  \texttt{\`{}shell\`{} is ignored when \`{}run\`{} is a map — shell is determined by
  map key resolution} \\
\hline
\texttt{run} map has unknown key & Error &
  \texttt{Unknown shell '\{key\}' in run map. Valid values: bash, sh, pwsh, cmd} \\
\hline
\texttt{shell} is explicit but not found on host & Plan-time Error &
  \texttt{Shell '\{shell\}' not found on this host (exec.LookPath returned empty)} \\
\hline
\texttt{run} map, no key resolves at run time & Run-time Error &
  \texttt{No shell from map \{keys\} available on this host} \\
\hline
\end{tabular}
\end{center}

% ─────────────────────────────────────────────────────────────────────────────
\subsection{Include Step Executor Contract}
\label{sec:include_executor_contract}
% ─────────────────────────────────────────────────────────────────────────────

The \texttt{include} step is the v2.0 inline composition primitive. It expands a referenced
runbook's steps directly into the caller's execution queue — no new run context, no isolated
variable scope, no separate audit entry.

\subsubsection*{Executor Algorithm}

When the executor reaches a \texttt{type: include} step it performs the following in order:

\begin{enumerate}
  \item \textbf{Resolve} the referenced runbook by relative path or registry ID.
  \item \textbf{Cycle check} — validate the include graph (see below). Fail with a hard error
        if any cycle is detected. This check must have already passed at load time; the
        executor re-asserts it as a defence-in-depth guard.
  \item \textbf{Evaluate \texttt{include.when}} — if the condition is present and evaluates to
        \texttt{false}, skip: emit no steps and advance past the include step.
  \item \textbf{Apply \texttt{include.with} overrides} to the current variable scope. These
        bindings are visible to all steps that follow (included or caller-defined).
  \item \textbf{Inline expand} — replace the include step in the execution queue with the
        referenced runbook's steps, each at \texttt{Depth+1} with their \texttt{Origin} field
        set to the referenced runbook path.
  \item \textbf{Continue execution} — the expanded steps execute as if they had been defined
        inline in the caller runbook.
\end{enumerate}

\subsubsection*{Key Architectural Properties}

\begin{itemize}
  \item \textbf{No new run context.} The included steps execute within the caller's run;
        \texttt{RunID}, variable scope, and audit log are all shared.
  \item \textbf{No separate audit entry for the include step itself.} Expanded steps appear
        in the caller's audit log each carrying an \texttt{Origin} breadcrumb of the form:
        \texttt{[included from: ./shared/db-preflight.runbook.yaml]}.
  \item \textbf{Shared variable scope.} Variable mutations made by included steps are
        visible to all subsequent steps in the caller, and vice versa.
  \item \textbf{Include step invisible in trace.} Only the expanded steps appear in the run
        trace. The include step itself is not emitted as a trace event.
\end{itemize}

\subsubsection*{Cycle Detection}
\label{sec:include_cycle_detection}

\paragraph{Load-time validation.}
When a runbook is loaded, the runtime builds a directed include graph over the full
transitive closure of all \texttt{type: include} references. A depth-first search (DFS)
with a visited set detects back-edges. A cycle found at any depth is a hard error:

\begin{verbatim}
cyclic include: A.yaml -> B.yaml -> A.yaml
\end{verbatim}

Cyclic include graphs are rejected before execution begins. This check covers both direct
and transitive includes.

\paragraph{Runtime note.}
Because cycles are validated at load time, the executor never needs to guard against
infinite include loops at execution time. The executor-level re-assertion in step 2 above
is a defence-in-depth check, not the primary enforcement point.

\paragraph{Future: Sub-Procedure Call.}
\texttt{type: include} is the v2.0 composition primitive: inline expansion into a shared
variable scope. A future \texttt{type: call} step type would provide isolated scope with
explicit input/output mapping, analogous to a function call. This pattern is deferred
post-v2.0; \texttt{type: include} covers the primary composition use case for the v2.0
release.

\subsubsection*{Step Post-flight (all step types)}
\begin{enumerate}
  \item Apply outcome evaluation: test \texttt{when:} expressions; record \texttt{OutcomeRecord}.
  \item Write captured variables to run state.
  \item Checkpoint run state to \texttt{.runbook/runs/<runID>/snapshots/<stepID>.json}.
  \item Write \texttt{step/completed} trace event (outcome, captures, elapsed time).
  \item Emit \texttt{StepCompleted} runtime event.
\end{enumerate}

\subsection{Phase 5: Run Completion}

When \texttt{Next} returns \texttt{io.EOF}:
\begin{enumerate}
  \item Write \texttt{run/completed} trace event (final outcome, step counts, elapsed time).
  \item Flush and sync the trace file.
  \item Compute SHA256 of the trace file; write \texttt{run/evidenceHash} record.
  \item Optionally archive the run directory (compress to \texttt{.runbook/archive/}).
  \item Close the event channel.
\end{enumerate}

% ─────────────────────────────────────────────────────────────────────────────
\section{Lifecycle Model}
% ─────────────────────────────────────────────────────────────────────────────

% Figure: Execution lifecycle state machines
\begin{figure}[htbp]
\centering
%% ── Run lifecycle ─────────────────────────────────────────────────────────────
\textit{\small Run lifecycle states:}\\[6pt]
\begin{tikzpicture}[node distance=0.8cm and 1.0cm, scale=0.9, transform shape]
  %% Main path
  \node[stepbox]                              (rs-created)   {created};
  \node[stepbox, right=0.8cm of rs-created]  (rs-planning)  {planning};
  \node[stepbox, right=0.8cm of rs-planning] (rs-ready)     {ready};
  \node[stepbox, right=0.8cm of rs-ready]    (rs-running)   {running};

  %% Terminal states (right of running)
  \node[stepbox, right=1.5cm of rs-running]         (rs-completed) {completed};
  \node[stepbox, below=0.8cm of rs-completed]       (rs-failed)    {failed};
  \node[stepbox, below=0.8cm of rs-failed]          (rs-cancelled) {cancelled};

  %% Suspension path (below running)
  \node[stepbox, below=1.3cm of rs-running]         (rs-suspended) {suspended};
  \node[stepbox, below=0.9cm of rs-suspended]       (rs-resuming)  {resuming};

  %% Arrows — main path
  \draw[gertarrow] (rs-created)   -- (rs-planning);
  \draw[gertarrow] (rs-planning)  -- (rs-ready);
  \draw[gertarrow] (rs-ready)     -- (rs-running);
  \draw[gertarrow] (rs-running)   -- (rs-completed);
  \draw[gertarrow] (rs-running)   -- (rs-failed);
  \draw[gertarrow] (rs-running)   to[out=-45, in=170] (rs-cancelled);

  %% Arrows — suspension cycle
  \draw[gertarrow] (rs-running)   -- (rs-suspended);
  \draw[gertarrow] (rs-suspended) -- (rs-resuming);
  \draw[gertarrow] (rs-resuming)  to[bend left=60] (rs-running);
\end{tikzpicture}

\medskip

%% ── Step lifecycle ────────────────────────────────────────────────────────────
\textit{\small Step lifecycle states:}\\[6pt]
\begin{tikzpicture}[node distance=0.8cm and 1.0cm, scale=0.9, transform shape]
  %% Main path
  \node[stepbox]                               (ss-pending)      {pending};
  \node[stepbox, right=1.2cm of ss-pending]    (ss-executing)    {running};
  \node[stepbox, right=1.8cm of ss-executing]  (ss-completed)    {completed};

  %% Waiting for input (below running, loops back)
  \node[stepbox, below=1.3cm of ss-executing]  (ss-waiting) {waiting};

  %% Terminal / error states (right of waiting)
  \node[stepbox, right=1.8cm of ss-waiting]    (ss-failed)       {failed};
  \node[stepbox, below=0.8cm of ss-failed]     (ss-skipped)      {skipped};
  \node[stepbox, below=0.8cm of ss-skipped]    (ss-denied)       {denied};
  \node[stepbox, below=0.8cm of ss-denied]     (ss-compensating) {compensating};

  %% Arrows
  \draw[gertarrow] (ss-pending)   -- (ss-executing);
  \draw[gertarrow] (ss-executing) -- (ss-completed);
  \draw[gertarrow] (ss-executing) -- (ss-waiting);
  \draw[gertarrow] (ss-waiting)   to[bend right=30] (ss-executing);
  \draw[gertarrow] (ss-executing) to[out=-30, in=165] (ss-failed);
  \draw[gertarrow] (ss-executing) to[out=-45, in=170] (ss-skipped);
  \draw[gertarrow] (ss-executing) to[out=-65, in=175] (ss-denied);
  \draw[gertarrow] (ss-executing) to[out=-75, in=175] (ss-compensating);
\end{tikzpicture}

\caption{Execution lifecycle --- run states (top) and step states (bottom).}
\label{fig:execution-lifecycle}
\end{figure}

\subsection{Run Start}

A run starts in one of two ways:

\begin{enumerate}
  \item \textbf{Direct invocation:} The CLI or test harness calls \texttt{Runtime.Start()}
        directly. This is the path for \texttt{gert exec}, \texttt{gert tui}, \texttt{gert debug},
        and \texttt{gert test}.
  \item \textbf{RPC invocation:} A client sends \texttt{exec/start} to the JSON-RPC server
        (\texttt{gert serve}). The server calls \texttt{Runtime.Start()} on behalf of the
        client and returns the \texttt{RunID} in the response. Subsequent step advancement
        is driven by \texttt{exec/next} RPC calls.
\end{enumerate}

\subsection{Step Advancement}

Step advancement is \emph{always client-driven}. The runtime does not auto-advance.
After \texttt{Start}, the client calls \texttt{Next} to execute each step in sequence.
This applies to all adapters, including headless CI mode. The only exception is the
\texttt{--auto} flag, which causes the CLI adapter to call \texttt{Next} in a tight loop,
producing the appearance of automatic advancement while retaining the single-step contract.

\subsection{Pause and Resume}

A run is \emph{paused} when \texttt{Next} returns a non-terminal result that requires external
input: interactive steps (\texttt{choice}, \texttt{decision}, \texttt{collector}) or approval
gates. The run remains paused until the required input is provided via \texttt{SubmitEvidence}
or \texttt{Approve}. The adapter is responsible for presenting the pause to the user.

The checkpoint written after each completed step is sufficient to resume a run after a
process restart. The checkpoint records: completed step IDs, current variable bindings,
current evidence values, and the trace sequence number. Resumption is initiated by calling
\texttt{Runtime.Resume(runID)} from any adapter.

% ─────────────────────────────────────────────────────────────────────────────
\subsection{Wait-for-Event Executor Contract}
\label{sec:wait_for_event_contract}
% ─────────────────────────────────────────────────────────────────────────────

\subsubsection*{WAITING Run State}

The \texttt{wait\_for\_event} step introduces a new \textbf{WAITING} run state that is distinct
from the existing user-initiated \textbf{PAUSED} state:

\begin{center}
\begin{tabular}{|l|l|p{7.5cm}|}
\hline
\textbf{State} & \textbf{Trigger} & \textbf{Exit Condition} \\
\hline
\texttt{RUNNING}  & \texttt{exec/start} or resume & Active step execution loop. \\
\hline
\texttt{PAUSED}   & User-initiated pause; interactive / approval step & \texttt{SubmitEvidence} or \texttt{Approve} RPC call from adapter. \\
\hline
\texttt{WAITING}  & \texttt{wait\_for\_event} step reached & Event Dispatcher delivers a matching event (\texttt{WAITING→RUNNING}); or timeout expires with \texttt{on\_timeout: fail} (\texttt{WAITING→FAILED}); or timeout expires with \texttt{on\_timeout: branch:<id>} (\texttt{WAITING→RUNNING}, jump to branch target). \\
\hline
\texttt{COMPLETED} & All steps finished & Terminal. \\
\hline
\texttt{FAILED}   & Unhandled error or \texttt{on\_timeout: fail} & Terminal (unless compensated in v2.1). \\
\hline
\end{tabular}
\end{center}

\texttt{WAITING} is a \emph{system-driven} suspension: the run is blocked on an external signal,
not on user input. Clients MUST differentiate the two states in UX (e.g.\ ``awaiting webhook''
vs.\ ``awaiting approval'') and MUST NOT send \texttt{exec/next} while the run is in
\texttt{WAITING} state.

\subsubsection*{Pause / Resume Protocol}

When the step execution loop reaches a \texttt{wait\_for\_event} step it MUST:

\begin{enumerate}
  \item \textbf{Serialise run state} — write the full run snapshot (step index, all variable
        bindings, call stack, active listener registration) to the run store (embedded BoltDB /
        SQLite, same store used for checkpoint snapshots).
  \item \textbf{Register event listener} — call
        \texttt{EventDispatcher.Register(runID, eventID, source, filter, payloadSchema, timeout)}.
        The Dispatcher returns a registration token.
  \item \textbf{Suspend} — set \texttt{run.state = WAITING}, write a
        \texttt{step/waiting} trace event (includes \texttt{event.source}, \texttt{event.id},
        timeout deadline), and return from the execution goroutine. The run goroutine is now
        idle; no \texttt{Next} call is outstanding.
\end{enumerate}

When the Event Dispatcher delivers a matching event to a \texttt{WAITING} run it MUST:

\begin{enumerate}
  \item Load the serialised run state from the run store via \texttt{store.Load(runID)}.
  \item Validate the payload against \texttt{step.event.payload\_schema} (JSON Schema Draft
        2020-12). Validation failure → \texttt{FAILED} with schema violation detail.
  \item Apply \texttt{step.event.filter} key-value predicates against the payload. Mismatch
        is treated as a non-matching event (listener remains registered; the event is
        discarded).
  \item Apply \texttt{step.capture} mappings: extract \texttt{json.<path>} fields from the
        payload and write them into \texttt{run.variables}.
  \item Write a \texttt{step/eventReceived} trace event (event source, sanitised payload
        summary, captured variable names).
  \item Set \texttt{run.state = RUNNING} and call \texttt{ExecuteNextStep(run)}.
\end{enumerate}

\subsubsection*{Executor Pseudocode}

\begin{minted}{bash}
// Suspend path
ExecuteStep(step: WaitForEvent, run: Run):
  store.Persist(run)                          // 1. durable checkpoint
  token = dispatcher.Register(               // 2. register listener
    run.id, step.event.id, step.event.source,
    step.event.filter, step.event.payload_schema,
    step.timeout)
  run.variables["gert.event."+step.event.id+".token"] = token  // HMAC token
  run.state = WAITING
  emit(step/waiting, {event_id: step.event.id, deadline: now()+step.timeout})
  return (suspended)                          // 3. goroutine exits

// Resume path
OnEventArrived(run_id, event_id, payload):
  run  = store.Load(run_id)                   // 1. rehydrate
  step = run.currentStep()                    //    (WaitForEvent step)
  ValidatePayload(payload, step.event.payload_schema)  // 2. schema check
  if not MatchFilter(payload, step.event.filter):      // 3. filter
    return  // discard; listener stays registered
  ApplyCapture(payload, step.capture, run.variables)   // 4. capture
  emit(step/eventReceived, {run_id, event_id})         // 5. trace
  run.state = RUNNING
  ExecuteNextStep(run)                        // 6. advance

// Timeout path
OnTimeout(run_id, event_id):
  run = store.Load(run_id)
  step = run.currentStep()
  switch step.on_timeout:
    case "fail":
      run.state = FAILED
      emit(step/failed, {reason: "wait_for_event timeout"})
    case "continue":
      run.state = RUNNING
      ExecuteNextStep(run)
    case "branch:<step-id>":
      run.state = RUNNING
      run.stepCursor = resolve(step-id)
      ExecuteNextStep(run)
\end{minted}

\subsubsection*{Run Persistence Requirement}

Runs in \texttt{WAITING} state MUST survive a \texttt{gert serve} process restart. The
serialised run snapshot written in step 1 above MUST include:

\begin{itemize}
  \item Current step index and depth.
  \item All variable bindings (including \texttt{gert.event.*} tokens).
  \item Call stack (for \texttt{include}-inlined steps).
  \item The active event listener registration (run ID, event ID, source, filter, timeout
        deadline as an absolute timestamp).
\end{itemize}

On \texttt{gert serve} restart, the server MUST re-register all persisted \texttt{WAITING}
listener records with the Event Dispatcher before accepting new connections. Listeners whose
timeout deadline has already passed are immediately moved to the timeout resolution path.

\subsubsection*{Serve-Only Constraint}

\texttt{wait\_for\_event} is only valid when running under \texttt{gert serve}. If a runbook
containing a \texttt{wait\_for\_event} step is executed via \texttt{gert run} (local CLI,
non-server mode), the Runtime Core MUST fail at dispatch time with:

\begin{verbatim}
step type 'wait_for_event' requires gert serve mode
\end{verbatim}

The Planner MAY emit a validation warning for \texttt{wait\_for\_event} steps in non-serve
context (aiding the user before execution begins), but the enforcement point is the Runtime
Core dispatch.

\subsubsection*{Security: HMAC Token for Webhook Transport}

When \texttt{event.source: webhook}, each \texttt{wait\_for\_event} registration generates a
cryptographically random one-time HMAC signing secret. The Event Dispatcher:

\begin{enumerate}
  \item Generates 32 random bytes via \texttt{crypto/rand}.
  \item Stores the secret alongside the listener registration in the run store.
  \item Exposes the HTTP endpoint \texttt{POST /events/\{run-id\}/\{event-id\}} and
        validates incoming requests using \texttt{HMAC-SHA256} over the raw request body
        (header: \texttt{X-Gert-Signature: sha256=<hex>}).
  \item Injects the secret into the run variable
        \texttt{gert.event.\{event-id\}.token} so the runbook author can communicate it to
        the external system (e.g.\ via a preceding \texttt{cli} step that calls
        \texttt{curl --data-urlencode}).
  \item Rejects requests with invalid or missing signatures with \texttt{HTTP 401}.
  \item Invalidates the secret after the first accepted delivery (one-time use).
\end{enumerate}

The token is intentionally injected as a run variable (not hard-coded in the runbook YAML)
so it is unique per run and never committed to source control.

% ─────────────────────────────────────────────────────────────────────────────
\subsection{Approve Step Executor Contract}
\label{sec:approve_executor_contract}
% ─────────────────────────────────────────────────────────────────────────────

The \texttt{approve} step is a first-class step type that pauses the run until one or more
authorised identities submit a sign-off decision via the Input Provider interface. Two new
capabilities — \textbf{business-day timeouts} (GAP-1) and \textbf{M-of-N quorum tracking}
(GAP-2) — extend the base approval gate without changing its fundamental PAUSED-state
semantics.

% ─────────────────────────────────────────────────────────────────
\subsubsection*{GAP-1: Business Calendar Engine}
\label{sec:business_calendar}
% ─────────────────────────────────────────────────────────────────

\paragraph{Motivation.}
Enterprise approval runbooks express deadlines in \emph{business days} (``respond within
3 business days''), not wall-clock durations. The existing \texttt{timeout} field is
wall-clock only and cannot satisfy this requirement.

\paragraph{Component responsibilities.}
The \textbf{Business Calendar Engine} is a stateless utility component within the Runtime
Core (no external process, no network calls for built-in calendars). Its responsibilities
are:

\begin{itemize}
  \item Determine whether a given instant is a business moment (Monday–Friday,
        non-holiday in the specified calendar and timezone).
  \item Calculate the wall-clock deadline that is exactly $N$ business days from a
        given start instant.
  \item Enumerate elapsed business days between two instants (for reporting and audit).
  \item Load and cache named business calendars; built-in calendars are compiled into
        the gert binary. Custom calendar definitions are deferred to v2.1.
\end{itemize}

\paragraph{Go interface.}

\begin{minted}{go}
type BusinessCalendar interface {
    // IsBusinessDay reports whether t falls on a business day (Mon–Fri, non-holiday)
    // in the given timezone.
    IsBusinessDay(t time.Time, tz *time.Location) bool

    // AddBusinessDays returns the instant that is exactly n business days after t
    // in the given timezone.  n must be >= 0.
    AddBusinessDays(t time.Time, days int, tz *time.Location) time.Time

    // ElapsedBusinessDays counts completed business days between start and end
    // in the given timezone.  Returns 0 if end <= start.
    ElapsedBusinessDays(start, end time.Time, tz *time.Location) int
}
\end{minted}

\paragraph{Built-in calendars for v2.0.}

\begin{center}
\begin{tabular}{|l|p{9cm}|}
\hline
\textbf{Calendar ID} & \textbf{Definition} \\
\hline
\texttt{default}    & Monday–Friday, no holidays. Appropriate for organisations that
                      define their own holiday schedule externally. \\
\hline
\texttt{us-federal} & Monday–Friday, US federal public holidays (New Year's Day,
                      Martin Luther King Jr.\ Day, Presidents' Day, Memorial Day,
                      Juneteenth, Independence Day, Labor Day, Columbus Day,
                      Veterans Day, Thanksgiving, Christmas). \\
\hline
\texttt{uk-banking} & Monday–Friday, UK bank holidays for England and Wales
                      (New Year's Day, Good Friday, Easter Monday, Early May Bank
                      Holiday, Spring Bank Holiday, Summer Bank Holiday, Christmas,
                      Boxing Day). \\
\hline
\end{tabular}
\end{center}

Custom calendar definitions (YAML-declared holiday lists referenced by a named ID) are
\textbf{deferred to v2.1} as a future-work item.

\paragraph{Timeout calculation at activation.}
The key architectural decision for business-day timeouts is a \textbf{one-time conversion
at step activation}: when the runtime dispatches an \texttt{approve} step that carries
\texttt{timeout\_business\_days}, it immediately calls:

\begin{verbatim}
deadline = calendar.AddBusinessDays(now, step.timeout_business_days, tz)
\end{verbatim}

The resulting \texttt{deadline} is a plain \texttt{time.Time} value stored in the run
state snapshot alongside the step. The existing timeout enforcement mechanism (the
Event Dispatcher's min-heap scheduler or the polling goroutine) then checks
\texttt{now >= deadline} at regular intervals — \emph{no business-calendar logic is
needed at check time}. This design means the Business Calendar Engine is invoked
exactly once per activate step and is entirely absent from the hot-polling path.

\begin{itemize}
  \item \textbf{timezone} is resolved from the step's \texttt{timezone} field (IANA
        timezone string, e.g.\ \texttt{America/New\_York}). If absent, the operator's
        local timezone is used; runbooks that require determinism SHOULD specify an
        explicit timezone.
  \item \textbf{calendar} is resolved from the step's \texttt{business\_calendar} field.
        Defaults to \texttt{default} if absent.
  \item If \texttt{timeout\_business\_days} and \texttt{timeout} (wall-clock) are both
        present on the same step, the \emph{earlier} of the two computed deadlines
        applies. A Planner warning is emitted to flag the redundancy.
\end{itemize}

% ─────────────────────────────────────────────────────────────────
\subsubsection*{GAP-2: Quorum Approval Tracker}
\label{sec:quorum_tracker}
% ─────────────────────────────────────────────────────────────────

\paragraph{Motivation.}
Financial and regulatory runbooks require ``3 of 5 department heads must approve''
semantics. The previous \texttt{approvals.min} field was a bare integer with no explicit
approver pool, making identity validation and audit reconstruction impossible.

\paragraph{Approval state per step instance.}
The runtime maintains an \texttt{ApprovalRecord} for each active \texttt{approve} step
instance:

\begin{minted}{go}
ApprovalRecord {
  step_id:     string
  run_id:      string
  mode:        all | any | quorum
  pool:        []string    // eligible approver role names
  required:    int         // 0 = all (for mode: all / any)
  approvals:   []Approval  // received approvals
  rejections:  []Approval  // received rejections
}

Approval {
  approver_id: string      // verified identity of the submitter
  role:        string      // which pool role the approver satisfied
  timestamp:   time.Time
  notes:       string      // optional free-text from the approver
}
\end{minted}

The \texttt{ApprovalRecord} is persisted to the run store at each update (same durable
store as run snapshots) so it survives a process restart.

\paragraph{Progression rules.}

\begin{center}
\begin{tabular}{|l|p{9cm}|}
\hline
\textbf{Mode} & \textbf{Proceed when} \\
\hline
\texttt{all}    & \texttt{len(approvals) == len(pool)} — every pool member has approved. \\
\hline
\texttt{any}    & \texttt{len(approvals) >= 1} — at least one pool member has approved. \\
\hline
\texttt{quorum} & \texttt{len(approvals) >= required} — the required count is reached. \\
\hline
\end{tabular}
\end{center}

Deduplication: each pool member may submit \textbf{exactly one} approval. If the same
\texttt{approver\_id} submits a second approval for the same step instance, the runtime
rejects the submission with \texttt{ErrDuplicateApproval} and returns an error to the
caller — it is not silently ignored.

Rejections are recorded in the audit trail and do \emph{not} block progression by
themselves, \emph{unless} quorum becomes mathematically impossible (see Deadlock
Detection below).

\paragraph{Deadlock detection.}
After \emph{every} approval or rejection event, the runtime evaluates:

\begin{minted}{text}
remaining_possible = len(pool) - len(rejections) - len(approvals)
if (len(approvals) + remaining_possible) < required:
    // quorum is impossible — fail immediately
    step.state = FAILED
    emit(step/failed, {reason: "QuorumImpossible", ...})
\end{minted}

For \texttt{mode: all}, \texttt{required} is treated as \texttt{len(pool)} for this
check. For \texttt{mode: any}, the check is skipped (one approval always suffices unless
the pool is empty, which is a schema validation error). If the check fires, the step
fails immediately with error code \texttt{QuorumImpossible}; the \texttt{on\_timeout}
policy does \emph{not} apply.

\paragraph{Audit trail.}
Every approval and rejection is appended to the run's append-only trace as a
\texttt{approve/decision} trace event containing:
\texttt{approver\_id}, \texttt{role}, \texttt{decision} (\texttt{approve}
or \texttt{reject}), \texttt{timestamp} (RFC 3339), and \texttt{notes}.

This satisfies electronic-signature audit requirements including SOC 2 Type II and
FDA 21 CFR Part 11.

\paragraph{Input Provider integration.}
Approval decisions are delivered via the Input Provider JSON-RPC interface. The
\texttt{approve} step executor registers a pending input request of type
\texttt{input/submitApproval} with the fields:

\begin{minted}{json}
{
  "run_id":   string,          // identifies the run
  "step_id":  string,          // identifies the approve step
  "decision": "approve"|"reject",
  "notes":    string           // optional
}
\end{minted}

Before recording the decision, the runtime:
\begin{enumerate}
  \item Resolves the caller's \texttt{approver\_id} from the authenticated session (or
        the \texttt{--as} flag in CLI mode).
  \item Verifies that \texttt{approver\_id} maps to at least one role in
        \texttt{approvals.pool}. Unknown identities receive \texttt{ErrNotAuthorised}.
  \item Checks for duplicate submission (\texttt{ErrDuplicateApproval} if already recorded).
  \item Appends the \texttt{Approval} or rejection record to the \texttt{ApprovalRecord}.
  \item Re-evaluates the progression rule and the deadlock-detection invariant.
  \item Persists the updated \texttt{ApprovalRecord} to the run store.
\end{enumerate}

% ─────────────────────────────────────────────────────────────────
\subsubsection*{Composition: Business-Day Timeout with Quorum Approval}
% ─────────────────────────────────────────────────────────────────

Both GAP-1 and GAP-2 \textbf{compose freely}. A step may declare
\texttt{timeout\_business\_days}, \texttt{timezone}, \texttt{business\_calendar}, and
\texttt{approvals} (\texttt{mode: quorum}) simultaneously. At step activation:

\begin{enumerate}
  \item The Business Calendar Engine computes the wall-clock \texttt{deadline}.
  \item The Quorum Tracker initialises an empty \texttt{ApprovalRecord}.
  \item The step enters PAUSED state; the timeout min-heap is armed with the deadline.
\end{enumerate}

The two mechanisms then run \textbf{independently}: the Quorum Tracker processes each
incoming approval/rejection without consulting the calendar, and the timeout scheduler
checks \texttt{now >= deadline} without consulting the quorum state. Whichever condition
fires first wins: quorum reached → step proceeds; deadline fired → \texttt{on\_timeout}
policy governs the outcome.

% ─────────────────────────────────────────────────────────────────────────────
\subsection{Display Step Executor Contract}
\label{sec:display_executor_contract}
% ─────────────────────────────────────────────────────────────────────────────

A \texttt{type: display} step renders template content to the operator's terminal
without blocking for input. It is the standard primitive for presenting collected
data, reports, and status summaries mid-flow.

\subsubsection*{Schema}

\begin{minted}[frame=lines,framesep=2mm,fontsize=\small]{yaml}
- step:
    id: show_health_report
    type: display
    title: "Service Health Report"     # step-list label
    display:
      content: "{{ .report }}"         # template body (required)
      format: text                     # text | markdown (default: text)
\end{minted}

\subsubsection*{Fields}

\begin{center}
\begin{tabular}{|l|l|l|p{7cm}|}
\hline
\textbf{Field} & \textbf{Type} & \textbf{Required} & \textbf{Description} \\
\hline
\texttt{content} & string & Yes & Template body rendered via \texttt{text/template}
against the current variable scope. YAML \texttt{|} and \texttt{>} blocks are supported.
Must not be empty. \\
\hline
\texttt{format} & enum & No & Render hint for the display surface. \texttt{text} (default):
pre-formatted plaintext, whitespace preserved. \texttt{markdown}: render with Markdown
formatting (bold, lists, headers, code blocks). \\
\hline
\end{tabular}
\end{center}

\subsubsection*{Execution Semantics}

The display executor:
\begin{enumerate}
  \item Evaluates \texttt{content} through the template engine (same evaluator used by
        \texttt{title}, \texttt{subtitle}, and \texttt{capture} expressions).
  \item Writes the rendered string to the operator's output stream.
        In CLI mode this is stdout; TUI clients may render into a styled panel.
  \item Returns immediately — the step never blocks for acknowledgment.
  \item Records a \texttt{step/completed} trace event. The rendered content is included
        in the event payload subject to the run's redaction policy.
  \item Never mutates the variable scope. \texttt{capture:} is disallowed on display steps
        (semantic validation rejects it).
\end{enumerate}

\subsubsection*{Surface Behaviour}

\begin{center}
\begin{tabular}{|l|p{5cm}|p{5cm}|}
\hline
\textbf{Surface} & \textbf{Behaviour} & \textbf{Notes} \\
\hline
CLI (TTY) & \texttt{fmt.Fprintln(stdout, rendered)} & ANSI passthrough if format is markdown \\
\hline
CLI (pipe / CI) & Same as TTY — output is emitted to stdout & Markdown stripped to plaintext \\
\hline
TUI & Rendered in the step panel / output card & \texttt{format: markdown} enables styled rendering \\
\hline
Dry-run & Emits trace event; no stdout write & Content is logged as-is \\
\hline
\end{tabular}
\end{center}

\subsubsection*{Disallowed Common Fields}

The following common \texttt{Step} fields are incompatible with \texttt{type: display}
and are rejected by semantic validation:

\begin{itemize}
  \item \texttt{capture:} — display produces no output values; nothing to bind.
  \item \texttt{retry:} — display cannot fail in a retriable sense.
  \item \texttt{contract:} — display has no side effects; a contract declaration is misleading.
\end{itemize}

\subsubsection*{Runbook Example}

\begin{minted}[frame=lines,framesep=2mm,fontsize=\small]{yaml}
flow:
  - iterate:
      id: collect_health_loop
      over: services
      as: svc
      steps:
        - step:
            id: check_svc
            type: include
            include:
              runbook: check-service.runbook.yaml
              with:
                service_host: "{{ .svc }}"
            capture:
              last_status: service_status

        - step:
            id: accumulate
            type: noop
            capture:
              report: "{{ .report }}{{ .svc }}: {{ .last_status }}\n"

  - step:
      id: show_report
      type: display
      title: "Health Report"
      display:
        content: |
          Service Health Report
          =====================

          {{ .report }}
        format: text

  - step:
      id: done
      type: end
      outcome:
        category: resolved
        code: health_report_complete
\end{minted}

The \texttt{type: display} step is the correct primitive for presenting accumulated
data to the operator. Using \texttt{type: noop} with a multiline \texttt{title:} is
explicitly an anti-pattern: \texttt{title} is a navigation label for the step list,
not a display surface.

% ─────────────────────────────────────────────────────────────────────────────
\subsection{Collector Field Validation Contract}
\label{sec:collector_field_validation}
% ─────────────────────────────────────────────────────────────────────────────

When the executor receives evidence from a \texttt{SubmitEvidence} call for a
\texttt{collector} step, it validates each field value against the field's declared type
and constraints \emph{before} writing any value to run state or the trace.

\subsubsection*{Validation Algorithm}

For each field in submission order:

\begin{enumerate}
  \item \textbf{Required check} — if the field is \texttt{required: true} and the
        submitted value is absent or empty, record a validation error.
  \item \textbf{Type coercion and format check} — attempt to parse or coerce the value
        to the declared type (see per-type rules below). A parse failure is a validation
        error.
  \item \textbf{Constraint check} — evaluate \texttt{validation.min},
        \texttt{validation.max}, \texttt{validation.step}, and option-list membership as
        applicable. A constraint violation is a validation error.
  \item \textbf{Re-prompt} — if any field produced a validation error, the executor sends
        a new \texttt{CollectorRequired} event back to the input provider with an
        \texttt{errors} map populated (field name → human-readable error message). The
        provider SHOULD re-render the form highlighting the invalid fields. The re-prompt
        counts as one retry attempt.
  \item \textbf{Max retries} — if validation still fails after \textbf{3} re-prompt
        attempts, the executor fails the step: \texttt{step.state = FAILED}, writes a
        \texttt{step/failed} trace event, and populates the \texttt{gert.error} reserved
        run variable with structured error detail (see below).
\end{enumerate}

\subsubsection*{Per-Type Validation and Storage Rules}

\begin{center}
\small
\begin{tabular}{|l|p{5.0cm}|p{3.6cm}|}
\hline
\textbf{Type} & \textbf{Validation} & \textbf{Variable Storage} \\
\hline
\texttt{text}
  & Required check; then optional length and pattern constraints:
    \texttt{validation.min\_length}, \texttt{validation.max\_length} (character
    count); \texttt{validation.pattern} (RE2 regex — see below);
    \texttt{validation.pattern\_hint} (shown on pattern failure).
    \texttt{multiline: true} is a UI hint only; no effect on validation or storage.
    If \texttt{format: yaml} or \texttt{format: json} is declared, structural
    parse validation is applied \emph{after} the pattern check (see below).
  & JSON string (raw; parsed form not stored) \\
\hline
\texttt{number}
  & Parse as IEEE~754 float64 (reject non-numeric input).
    Check \texttt{validation.min} and \texttt{validation.max} (both inclusive).
    Check \texttt{validation.step}: \texttt{(value - min) \% step < epsilon}.
  & JSON number (\texttt{float64}) \\
\hline
\texttt{integer}
  & Parse as int64; reject fractional input.
    Same \texttt{min}/\texttt{max}/\texttt{step} checks as \texttt{number}.
  & JSON number (\texttt{int64}) \\
\hline
\texttt{date}
  & Must match \texttt{YYYY-MM-DD} (RFC~3339 full-date).
    Validate \texttt{validation.min}/\texttt{max} as date strings (lexicographic
    comparison is correct for ISO dates).
    Wall clock at submission time; timezone UTC.
  & JSON string (\texttt{YYYY-MM-DD}) \\
\hline
\texttt{datetime}
  & Must be valid RFC~3339 (ISO~8601 with time offset).
    Validate \texttt{validation.min}/\texttt{max} as RFC~3339 strings; normalise to
    UTC before comparison.
    Wall clock at submission time.
  & JSON string (RFC~3339, e.g.\ \texttt{2026-04-18T14:00:00Z}) \\
\hline
\texttt{boolean}
  & Truthy/falsy coercion:
    \texttt{true}, \texttt{"true"}, \texttt{1}, \texttt{"yes"} → \texttt{true};
    \texttt{false}, \texttt{"false"}, \texttt{0}, \texttt{"no"} → \texttt{false}.
    Any other value is a validation error.  No min/max constraints.
  & JSON boolean \\
\hline
\texttt{select}
(\texttt{multiple: false})
  & Submitted value must exactly match one of the declared
    \texttt{options[*].value} strings.  If \texttt{options\_from} was used,
    validation is against the dynamically fetched list (§\ref{sec:dynamic-options-protocol}).
  & JSON string \\
\hline
\texttt{select}
(\texttt{multiple: true})
  & All submitted values must be members of the options list.
    Validate \texttt{validation.min\_selections} and
    \texttt{validation.max\_selections} if declared.
    Duplicate values within the submission array are a validation error.
  & JSON array of strings \\
\hline
\texttt{choice}
(\texttt{multiple: true})
  & Same rules as \texttt{select} with \texttt{multiple: true}, applied to the
    \texttt{choice} step type.
  & JSON array of strings \\
\hline
\texttt{file},
\texttt{image}
  & File size and MIME type constraints (see §\ref{subsec:collector-contract}).
    SHA-256 computed by executor after artifact is stored.
  & Artifact reference (evidence store; not a plain variable) \\
\hline
\texttt{url}
  & Must parse as a valid URL (scheme required; HTTP and HTTPS accepted by
    default; custom schemes allowed).
  & JSON string \\
\hline
\end{tabular}
\end{center}

\subsubsection*{\texttt{multiline: true} Hint}

The \texttt{multiline: true} flag on a \texttt{text} field is a \textbf{UI hint only}.
The executor treats the submitted value as a plain string regardless of this flag.
Input providers SHOULD render a multi-line textarea (or open \texttt{\$EDITOR}) rather
than a single-line input box when \texttt{multiline: true} is set.

% ─────────────────────────────────────────────────────────────────────────────
\subsubsection*{\texttt{text} Field Pattern Validation}
% ─────────────────────────────────────────────────────────────────────────────

\texttt{text} fields support optional length, pattern, and structural-format
constraints applied \emph{after} the required check.  The full sequence for a
\texttt{text} field is:

\begin{enumerate}
  \item \textbf{Required check} — same as all field types (step 1 of the general
        algorithm above).
  \item \textbf{Min-length check} — if \texttt{validation.min\_length} is declared,
        fail if \texttt{len(value) < min\_length} (character count, not bytes).
  \item \textbf{Max-length check} — if \texttt{validation.max\_length} is declared,
        fail if \texttt{len(value) > max\_length}.
  \item \textbf{Pattern check} — if \texttt{validation.pattern} is declared, compile
        and test the full value against the RE2 regex.  On failure, the re-prompt
        error message is \texttt{validation.pattern\_hint} if present; otherwise the
        generic message \emph{``Value does not match required format''} is used.
  \item \textbf{Format check} — if \texttt{format: yaml} is declared, attempt
        \texttt{yaml.Unmarshal([]byte(value), \&interface\{\})}.  On error, re-prompt
        with \emph{``Value must be valid YAML.''}.  If \texttt{format: json} is
        declared, attempt \texttt{json.Unmarshal([]byte(value), \&interface\{\})}.
        On error, re-prompt with \emph{``Value must be valid JSON.''}.  Both checks
        are \textbf{parse-only}: the parsed result is discarded; the value remains a
        raw string in the run variable.  Downstream steps use the \texttt{fromYAML}
        and \texttt{fromJSON} template functions for traversal
        (see §\ref{subsec:expressions}).  If both \texttt{validation.pattern}
        and \texttt{format} are declared, both must pass; the pattern check (step 4)
        runs first.
  \item \textbf{Re-prompt} — any failure at steps 2--5 triggers a re-prompt.  The
        3-attempt limit is a \textbf{combined total} across all validation steps: if
        the operator fails step 2 once, step 4 once, and step 5 once, the step fails
        on the third attempt regardless of which check failed each time.
\end{enumerate}

\paragraph{RE2 vs ECMA~262.}
The gert schema specification uses the term \emph{ECMA~262 regex} for compatibility
with JSON Schema tooling and browser-based UI validators.  At runtime, gert uses Go's
\texttt{regexp} package, which implements RE2 semantics — not a full ECMA~262 engine.

Patterns \textbf{must} be valid RE2 expressions.  ECMA~262 features absent from RE2
(lookahead, lookbehind, backreferences) are \textbf{rejected at schema load time} with
a validation error naming the offending field:

\begin{minted}{yaml}
error: step "gather-details", field "ticket_ref":
  validation.pattern uses lookahead (?=...) which is not supported in RE2;
  replace with an anchored RE2 equivalent
\end{minted}

Runbook authors targeting both the VS~Code extension (which uses a JavaScript regex
engine) and the gert runtime (RE2) should restrict patterns to the safe common subset:
character classes, quantifiers, anchors, alternation, and non-capturing groups.

% ─────────────────────────────────────────────────────────────────────────────
\subsubsection*{Conditional Field Evaluation Contract}
\label{sec:conditional-field-eval}
% ─────────────────────────────────────────────────────────────────────────────

\texttt{collector} fields MAY declare a \texttt{when} expression.  The executor
evaluates \texttt{when} before rendering each field, determining whether the field
is shown, prompted, and bound to a variable.

\paragraph{Evaluation model.}
Fields are evaluated left-to-right in array order:

\begin{enumerate}
  \item Before rendering field $i$, the executor evaluates its \texttt{when}
        expression against:
        \begin{itemize}
          \item \textbf{Collected values so far} — the values returned by the input
                provider for fields $0 \ldots i-1$ in the current step (those fields
                whose \texttt{when} evaluated to \texttt{true} and were rendered).
          \item \textbf{Global variable scope} — the run's accumulated variable
                bindings from all previously completed steps.
        \end{itemize}
  \item If \texttt{when} evaluates to \textbf{false}: the field is \textbf{not
        rendered, not prompted, and its variable binding is skipped entirely}.  The
        field's \texttt{id} is NOT added to the step's variable scope for this run.
  \item If \texttt{when} evaluates to \textbf{true}, or if \texttt{when} is absent
        (unconditionally visible), the field is rendered and validated normally.
\end{enumerate}

\noindent
This is a \textbf{pull model}: the executor drives field-by-field evaluation in order;
the input provider renders one field at a time, or a full form with some fields
conditionally hidden — both are valid provider strategies.

\paragraph{Input provider integration.}
Two delivery strategies are both valid; providers declare their preference via
capability flags (see §\ref{subsec:collector-contract}):

\begin{itemize}
  \item \textbf{Full-form delivery} (recommended for rich UI providers): the executor
        sends the complete \texttt{fields} array in a single \texttt{provider/collect}
        request.  Each field carries a pre-evaluated \texttt{visible: bool} flag
        reflecting the current \texttt{when} result at dispatch time.  Rich UI
        providers SHOULD implement dynamic client-side show/hide as the user fills in
        earlier fields, re-evaluating \texttt{when} locally against live form state so
        the form responds without a round-trip to the executor.
  \item \textbf{Streaming delivery} (recommended for CLI providers): the executor
        sends only the currently-visible fields, evaluating each \texttt{when}
        expression in real time as previous field values are received.  This is the
        natural mode for the built-in \texttt{prompt} provider, which renders fields
        sequentially at the terminal.  Fields where \texttt{visible: false} are
        simply never sent.
\end{itemize}

\paragraph{Variable binding on skip.}
When a field is skipped (\texttt{when} = \texttt{false}):

\begin{itemize}
  \item Its \texttt{id} is NOT added to the step's variable bindings.
  \item Template expressions in subsequent steps that reference the skipped field's
        variable resolve to the \emph{zero value}: \texttt{""} for string,
        \texttt{0} for number, \texttt{false} for boolean.
  \item If the skipped variable is referenced with \texttt{.require} in a downstream
        template expression, the executor produces a template error at that step.
  \item Runbook authors should guard downstream template expressions that may depend
        on conditionally-bound fields:
\end{itemize}

\begin{verbatim}
{{/* Safe: test for empty before using a conditionally-bound variable */}}
{{ if .ticket_ref }}JIRA: {{ .ticket_ref }}{{ else }}(no ticket){{ end }}
\end{verbatim}

\paragraph{Load-time validation — forward reference check.}
At schema load time, the parser performs a forward-reference check for every field
that declares \texttt{when}:

\begin{enumerate}
  \item Parse the \texttt{when} expression and extract all variable references.
  \item For each referenced variable, determine whether it is bound by a field
        declared \emph{earlier} in the same \texttt{fields} array (by array index).
  \item If any referenced variable is bound by a field declared \emph{later} in the
        same \texttt{fields} array, \textbf{reject the runbook at load time} with a
        validation error:
\end{enumerate}

\begin{minted}{yaml}
error: step "gather-details", field "region" (index 2):
  when expression references "country" which is bound by field at index 4;
  forward references in when expressions are not allowed
\end{minted}

\noindent
Variables from the global scope (prior steps) and the run context are always valid
references in \texttt{when} expressions; only same-step forward references are
rejected.

\subsubsection*{Variable Storage — Downstream Expression Semantics}

Typed storage ensures that template comparison operators in downstream steps work without
manual coercion:

\begin{verbatim}
{{/* number/integer: .amount is float64, not a string */}}
{{ if gt .amount 50000.0 }}HIGH_VALUE{{ end }}

{{/* boolean: .confirmed is bool, not "yes"/"no" */}}
{{ if .confirmed }}PROCEED{{ end }}

{{/* select/multiple or choice/multiple: .systems is []string */}}
{{ if has .systems "payments-api" }}CRITICAL{{ end }}

{{/* date/datetime: .start_date is ISO 8601 string; parseable by time.Parse */}}
{{ if lt .start_date "2026-07-01" }}BEFORE_CUTOVER{{ end }}
\end{verbatim}

\subsubsection*{Error Variable Schema}

When a \texttt{collector} step exhausts its re-prompt attempts and fails, the executor
populates the \texttt{gert.error} reserved run variable:

\begin{minted}{json}
{
  "step_id":  "gather-approval-details",
  "kind":     "collector.validation_failed",
  "attempts": 3,
  "fields": {
    "amount":     "value 49000 is below minimum 50000",
    "start_date": "2025-12-01 is before minimum 2026-01-01",
    "department": "\"legal\" is not a valid option"
  }
}
\end{minted}

The \texttt{gert.error} variable is a reserved JSON object populated on any step
failure. A subsequent \texttt{branch} step may inspect it to route to an error-handling
path. Each new step failure overwrites the previous value.

% ─────────────────────────────────────────────────────────────────────────────
\subsubsection*{Ephemeral Field Handling}
\label{sec:ephemeral-field-handling}
% ─────────────────────────────────────────────────────────────────────────────

A \texttt{collector} field may declare \texttt{ephemeral: true} to indicate that its
value is sensitive and must never be written to any durable store.

\paragraph{Audit redaction.}
After a field value passes all validation checks, the executor writes the step's output
to the JSONL trace file.  For fields where \texttt{ephemeral: true}:

\begin{itemize}
  \item The JSONL audit record stores the string literal \texttt{"[REDACTED]"} as the
        field value — not the actual value.
  \item The \emph{actual} value IS placed in the in-memory \texttt{VariableScope} for
        the duration of the run and is available to downstream steps via normal template
        expressions.
  \item The actual value is \textbf{NEVER written to disk} in any log, trace, snapshot,
        or checkpoint file.  This includes the run trace (\texttt{trace.jsonl}), the
        per-step checkpoint (\texttt{snapshots/<stepID>.json}), and the evidence hash
        record.
\end{itemize}

\paragraph{Go interface contract.}
The in-memory \texttt{VariableScope} type MUST expose the method:

\begin{verbatim}
IsEphemeral(key string) bool
\end{verbatim}

The JSONL trace writer calls \texttt{IsEphemeral} for every variable before serialising
a \texttt{step/completed} or \texttt{collector/submitted} event.  If the method returns
\texttt{true}, the writer substitutes \texttt{"[REDACTED]"} for the value in the emitted
JSON.  No other layer in the stack applies this substitution — the redaction responsibility
belongs exclusively to the trace writer.

\paragraph{Tool step compatibility.}
Ephemeral values MAY be passed as arguments to \texttt{type: tool} steps.  The tool
receives the plaintext value.  This is intentional: \texttt{ephemeral: true} is an
\emph{audit redaction} feature, not end-to-end encryption.  Runbook authors are
responsible for using ephemeral fields only with tools that handle the value securely.

\paragraph{Suspension/resume constraint.}
When a run is suspended (e.g., a \texttt{wait\_for\_event} or \texttt{approve} step
puts the run into \texttt{WAITING} state), the in-memory \texttt{VariableScope} is not
persisted.  Because ephemeral values exist only in memory, they are \textbf{dropped} on
suspension.  On resume, any downstream step that requires an ephemeral variable will
find it absent from the reconstructed scope.

\begin{tcolorbox}[colback=yellow!10,colframe=orange!60,title=Constraint: ephemeral variables and suspension]
Steps that depend on ephemeral variables MUST appear before any \texttt{wait\_for\_event}
or \texttt{approve} step that could suspend the run.  The schema validator SHOULD emit a
load-time warning (not a hard error in v2.0) if an ephemeral variable is referenced in a
step that follows a suspendable step in the same runbook's execution path.
\end{tcolorbox}

\subsection{Cancellation}

\texttt{RunHandle.Cancel(ctx, reason)} initiates graceful cancellation:
\begin{enumerate}
  \item The currently-executing subprocess (if any) receives \texttt{SIGTERM} followed by
        \texttt{SIGKILL} after a 5-second grace period.
  \item The runtime writes a \texttt{run/cancelled} trace event with the reason and actor.
  \item Subsequent \texttt{Next} calls return \texttt{ErrRunCancelled}.
  \item The run state is checkpointed as \texttt{cancelled} (resumable).
\end{enumerate}

The cancellation contract does not invoke compensation handlers in v2.0. Saga/compensation
(stepping backward through a compensation chain on failure) is deferred to v2.1 as specified
in §09 (Open Questions).

\subsection{Crash Recovery}

If the gert process crashes mid-run, the trace and checkpoint files remain on disk in a
consistent state due to the append-only, sync-after-write trace discipline. On restart,
the user (or the VS Code extension) can call \texttt{Runtime.Resume(runID)} to continue
from the last completed step. The runtime reads the checkpoint to reconstruct run state.
Steps already recorded as completed in the trace are not re-executed.

% ─────────────────────────────────────────────────────────────────────────────
\section{Concurrency Model}
% ─────────────────────────────────────────────────────────────────────────────

\subsection{Per-Process Run Isolation}

The v2 runtime is \textbf{single-run-per-process}. Each \texttt{gert exec} invocation is a
separate OS process. \texttt{gert serve} is the exception: it is a long-running process that
may host multiple sequential runs (one at a time) for a single VS Code extension session.
Concurrent runs within a single \texttt{gert serve} instance are explicitly out of scope
for v2.0 (see §09).

\paragraph{Rationale.} Governance enforcement, file-level trace writes, and subprocess management
are dramatically simpler with one run per process. Process isolation also provides a natural
crash boundary: a failed run cannot corrupt the state of a concurrent run.

\subsection{Goroutine Model}

Within a single run, the concurrency model is:

\begin{itemize}
  \item \textbf{One goroutine per run:} the step execution loop runs on a single goroutine.
        Steps execute sequentially; there is no goroutine-per-step parallelism in v2.0.
  \item \textbf{Event fan-out goroutine:} a second goroutine reads from an internal event
        buffer and writes to the \texttt{Events()} channel. This decouples event emission
        (which must be synchronous with step execution) from event consumption (which may
        be slow, e.g., a WebSocket write).
  \item \textbf{Subprocess goroutine:} when a \texttt{cli} step spawns a subprocess, a
        goroutine is used to drain stdout/stderr concurrently, preventing deadlock on
        full pipe buffers.
  \item \textbf{Tool RPC goroutine:} tool invocations over JSON-RPC or MCP use a goroutine
        per call, managed by the Tool Runtime's connection pool.
\end{itemize}

\subsection{Thread Safety at Component Boundaries}

\begin{itemize}
  \item \texttt{RunHandle} is \textbf{not safe for concurrent use}. The adapter must
        serialize all calls. \texttt{gert serve} enforces this with a per-session mutex.
  \item The \texttt{Events()} channel is safe for a single consumer goroutine.
  \item The \texttt{GovernanceEngine} is \textbf{read-only after construction} and therefore
        safe for concurrent read access (relevant if future versions support parallel steps).
  \item The Trace Writer is protected by an internal mutex; write calls are safe to issue
        from any goroutine within the run.
\end{itemize}

% ─────────────────────────────────────────────────────────────────────────────
\section{\texttt{gert serve} Integration}
% ─────────────────────────────────────────────────────────────────────────────

\texttt{gert serve} is the JSON-RPC 2.0 server that powers the VS Code extension, the MCP
tools layer, and any future web client. In the v2 architecture it is classified as an
\textbf{adapter in the API/Adapter Layer}.

\subsection{Architectural Position}

\texttt{gert serve} is a process that:
\begin{enumerate}
  \item Starts a long-lived stdio JSON-RPC 2.0 listener (newline-delimited JSON over stdin/stdout).
  \item On receiving \texttt{exec/start}, calls \texttt{Runtime.Start()} and stores the
        resulting \texttt{RunHandle} in a session map (keyed by \texttt{RunID}).
  \item On receiving \texttt{exec/next}, calls \texttt{RunHandle.Next()} and returns the
        \texttt{StepResult} as the RPC response.
  \item Drains the \texttt{Events()} channel in a background goroutine and forwards each
        event as a JSON-RPC notification to the client.
  \item On receiving \texttt{exec/approve}, \texttt{exec/submitEvidence}, or
        \texttt{exec/cancel}, dispatches to the corresponding \texttt{RunHandle} method.
\end{enumerate}

\subsection{Event Forwarding}

Every event emitted by the Runtime Core is forwarded to the client as a JSON-RPC
\emph{notification} (no \texttt{id} field) with method \texttt{event/\textit{eventType}}.
The payload is the event's JSON representation. The VS Code extension subscribes to
these notifications to update its UI without polling.

Events forwarded over WebSocket (for a future web adapter) follow the same schema;
only the transport changes.

\subsection{Crash Recovery in \texttt{gert serve}}

If the VS Code extension reconnects after a crash (e.g., the gert process was killed),
\texttt{gert serve} exposes \texttt{exec/list} and \texttt{exec/resume} methods. The
extension discovers the interrupted run, calls \texttt{exec/resume}, and receives a new
\texttt{RunHandle} over the same session. The trace and checkpoint on disk ensure no
steps are lost.

% ─────────────────────────────────────────────────────────────────────────────
\subsection{Event Dispatcher}
\label{sec:event_dispatcher}
% ─────────────────────────────────────────────────────────────────────────────

The \textbf{Event Dispatcher} is a new runtime component that lives exclusively inside the
\texttt{gert serve} process. It is not instantiated by \texttt{gert run}.

\subsubsection*{Responsibilities}

\begin{itemize}
  \item \textbf{Listener lifecycle} — maintain an in-memory registry of active event listeners
        (indexed by \texttt{run\_id + event\_id}). Each entry records: source type, filter
        predicates, payload schema URL, timeout deadline, and HMAC secret (webhook source only).
        Persist the registry to the run store so it survives process restart.
  \item \textbf{Webhook transport} — expose \texttt{POST /events/\{run-id\}/\{event-id\}}
        authenticated with \texttt{HMAC-SHA256} (see security note in
        §\ref{sec:wait_for_event_contract}). Parse the request body as JSON and dispatch to
        the matching listener. Return \texttt{HTTP 202 Accepted} on success; \texttt{404} if
        no listener is registered; \texttt{401} on signature failure; \texttt{410 Gone} if
        the run is no longer in \texttt{WAITING} state.
  \item \textbf{Channel transport} — maintain an in-process \texttt{map[string]chan Payload}
        keyed by \texttt{event.id}. Any component within the same \texttt{gert serve} process
        can post to a named channel via \texttt{EventDispatcher.Send(eventID, payload)}.
  \item \textbf{Message broker transport} — delegate to a pluggable \texttt{MessageBroker}
        interface (implementation-defined for v2.0; a no-op stub ships by default). The
        interface follows the same pattern as the \texttt{ToolTransport} plug-in model: a
        registered factory function keyed by broker type string.
  \item \textbf{Signal transport} — register \texttt{os/signal} handlers scoped to a specific
        run. Signal handling is multiplexed: the same OS signal may be registered by multiple
        active listeners (e.g.\ \texttt{SIGUSR1} on two concurrent waiting runs). Each
        listener receives its own notification; signal handlers are deregistered on listener
        expiry or run completion.
  \item \textbf{Timeout enforcement} — maintain a min-heap of timeout deadlines. A background
        goroutine fires on the earliest deadline and invokes \texttt{OnTimeout(runID, eventID)}.
\end{itemize}

\subsubsection*{Go Interface}

\begin{minted}{go}
type EventDispatcher interface {
    // Register creates a listener for an inbound event on a waiting run.
    // Returns a one-time HMAC token (non-empty for webhook source only).
    Register(ctx context.Context, reg ListenerRegistration) (token string, err error)

    // Deregister removes a listener (called on run completion or cancellation).
    Deregister(runID, eventID string) error

    // Send posts a payload to a named in-process channel listener.
    Send(eventID string, payload json.RawMessage) error

    // RestoreFromStore re-registers all WAITING listeners after process restart.
    RestoreFromStore(store RunStore) error
}

type ListenerRegistration struct {
    RunID          string
    EventID        string
    Source         EventSource          // webhook | message | signal | channel
    Filter         map[string]string
    PayloadSchema  string               // JSON Schema URL, may be empty
    TimeoutAt      time.Time
    OnTimeout      TimeoutPolicy        // fail | continue | branch:<step-id>
}
\end{minted}

\subsubsection*{Availability Constraint}

The Event Dispatcher is instantiated by \texttt{gert serve} only. It is \textbf{not} part of
the \texttt{Runtime} interface exposed to step executors directly; instead, the Runtime Core
receives a \texttt{DispatcherHandle} at construction time (nil in \texttt{gert run} mode).
The \texttt{wait\_for\_event} dispatch path checks for a nil handle and fails immediately
with the serve-mode error before attempting any state serialisation.

\subsection{RPC Method Summary}

\begin{minted}{bash}
exec/start          → Start a new run (returns RunID, first StepResult)
exec/next           → Advance run by one step (returns StepResult)
exec/approve        → Submit approval decision for current gate
exec/submitEvidence → Submit evidence for interactive steps (choice/decision/collector)
exec/cancel         → Cancel the active run
exec/list           → List all runs for this workspace (active + completed)
exec/resume         → Resume a previously interrupted run
runbook/validate    → Validate a runbook YAML (parse + plan, no exec)
runbook/diagram     → Generate a Mermaid or ASCII diagram
event/<type>        → Notification from server to client (not a request)

// HTTP endpoints (gert serve --http, Event Dispatcher)
POST /events/{run-id}/{event-id}  → Deliver a webhook event to a waiting run
                                    (HMAC-SHA256 authenticated)
\end{minted}
